\documentclass[12pt,a4paper]{article}
\usepackage[utf8]{inputenc}
\usepackage[spanish]{babel}
\usepackage{geometry}
\usepackage{hyperref}
\usepackage{listings}
\usepackage{xcolor}
\usepackage{longtable}
\usepackage{enumitem}
\usepackage{titlesec}
\usepackage{fancyhdr}
\usepackage{graphicx}

\geometry{margin=2.5cm}
\setlength{\parindent}{0pt}
\setlength{\parskip}{6pt}

\lstset{
  basicstyle=\ttfamily\small,
  backgroundcolor=\color{gray!10},
  frame=single,
  breaklines=true,
  numbers=left,
  numberstyle=\tiny,
  language=bash,
  showstringspaces=false
}

\pagestyle{fancy}
\fancyhf{}
\fancyhead[L]{\small\textit{Pixel Store --- Guía de Despliegue}}
\fancyhead[R]{\thepage}
\fancyfoot[C]{\small Documento interno --- Versión 1.0}

\title{\textbf{Guía de Despliegue\\Pixel Store}}
\author{Documentación Técnica}
\date{\today}

\begin{document}
\maketitle
\thispagestyle{empty}
\newpage

\tableofcontents
\newpage

\section{Descripción General}

Pixel Store es una tienda digital dividida en dos proyectos principales:

\begin{itemize}
  \item \textbf{Frontend (React + Vite + TypeScript)} --- Aplicación web SPA en \texttt{src/}.
  \item \textbf{Backend (Node.js + Express + Socket.io)} --- API REST + WebSockets en \texttt{server/}.
  \item \textbf{Store Estático} --- Versión alternativa en \texttt{C:/items/peek-store-static/}.
\end{itemize}

\subsection{Servidor de Producción}

\begin{lstlisting}
Host:   177.7.41.162
User:   root
Pass:   Pixelstore34.5
OS:     Linux
\end{lstlisting}

\subsection{Dominios}

\begin{lstlisting}
Frontend (React SPA):   https://vitria.cc  (Vercel, push a main)
Backend (API):          https://api.vitria.cc  (VPS Nginx + PM2)
Static store:           (en el VPS, /var/www/frontend)
\end{lstlisting}

\newpage

\section{Estructura del Proyecto}

\subsection{Raíz del Frontend}

\begin{lstlisting}
pixel_store_home_page_58ruox_dualiteproject/
  src/
    App.tsx                 -- Router principal
    main.tsx                -- Entry point
    index.css               -- Estilos globales
    services/api.ts         -- API client + Socket.io client
    pages/                  -- Paginas (12)
      Account.tsx, Admin.tsx, Catalog.tsx, Chat.tsx,
      Checkout.tsx, Fortnite.tsx, GameItems.tsx,
      Groups.tsx, Home.tsx, OrderDetails.tsx,
      Reviews.tsx, RobuxCatalog.tsx
    components/
      admin/                -- Componentes del panel admin
        ChatsTab.tsx, OrdersTab.tsx, Mm2Tab.tsx,
        LimitedsTab.tsx, DeliveriesTab.tsx,
        FortniteTab.tsx, ProductsTab.tsx, etc.
      fortnite/             -- Componentes de Fortnite
      Navbar.tsx, Footer.tsx, footer2.tsx, etc.
  server/
    index.js                -- Entry point del backend (Express + Socket.io)
    database.js             -- Inicializacion de bases de datos JSON (lowdb)
    routes/
      admin.js, auth.js, chat.js, coupons.js,
      fortnite.js, orders.js, products.js,
      reviews.js, roblox.js, robloxGamepasses.js,
      robloxGroups.js, robloxPlaces.js
    services/
      discordService.js     -- Embed builder para Discord
    db/                     -- Archivos JSON (datos persistidos)
    uploads/                -- Archivos subidos
  dist/                     -- Build de produccion del frontend
  vite.config.ts            -- Configuracion de Vite
  package.json              -- Dependencias del frontend
  vercel.json               -- Configuracion de despliegue Vercel
\end{lstlisting}

\subsection{Rutas del Backend (API)}

\begin{longtable}{|l|l|l|}
\hline
\textbf{Método} & \textbf{Ruta} & \textbf{Descripción} \\ \hline
\endfirsthead
\hline
\textbf{Método} & \textbf{Ruta} & \textbf{Descripción} \\ \hline
\endhead
GET  & /api/products & Listar productos \\ \hline
POST & /api/products & Crear producto (admin) \\ \hline
GET  & /api/orders & Listar pedidos (filtro por usuario o admin) \\ \hline
POST & /api/orders & Crear pedido \\ \hline
GET  & /api/orders/:id & Detalle de pedido \\ \hline
PUT  & /api/orders/:id & Actualizar pedido (admin) \\ \hline
POST & /api/chat/message & Enviar mensaje / crear chat \\ \hline
GET  & /api/chat/:id/messages & Obtener mensajes de un chat \\ \hline
GET  & /api/chat/unread-count & Contador de no leídos \\ \hline
POST & /api/auth/login & Login de admin \\ \hline
POST & /api/reviews & Crear reseña \\ \hline
GET  & /api/reviews & Listar reseñas \\ \hline
GET  & /api/admin/stats & Estadísticas del panel admin \\ \hline
\end{longtable}

\subsection{Socket.io --- Eventos en Tiempo Real}

\begin{longtable}{|l|l|l|}
\hline
\textbf{Evento} & \textbf{Dirección} & \textbf{Propósito} \\ \hline
\endfirsthead
\hline
\textbf{Evento} & \textbf{Dirección} & \textbf{Propósito} \\ \hline
\endhead
\texttt{join-chat} & cliente $\rightarrow$ servidor & Unirse a sala de chat \\ \hline
\texttt{join-order} & cliente $\rightarrow$ servidor & Unirse a sala de orden \\ \hline
\texttt{join-admin} & cliente $\rightarrow$ servidor & Unirse a sala global de admin \\ \hline
\texttt{new-message} & servidor $\rightarrow$ cliente & Mensaje nuevo en chat \\ \hline
\texttt{new-chat} & servidor $\rightarrow$ cliente & Chat nuevo creado \\ \hline
\texttt{update-chat-list} & servidor $\rightarrow$ cliente & Actualizar lista de chats \\ \hline
\texttt{notification-admin} & servidor $\rightarrow$ cliente & Notificación para admin \\ \hline
\texttt{notification-\{userId\}} & servidor $\rightarrow$ cliente & Notificación para usuario \\ \hline
\texttt{order-updated} & servidor $\rightarrow$ cliente & Orden actualizada \\ \hline
\texttt{chat-ended} & servidor $\rightarrow$ cliente & Chat finalizado \\ \hline
\end{longtable}

\newpage

\section{Despliegue del Backend}

\subsection{Subir Archivos al VPS}

Usar \texttt{pscp} (Windows) o \texttt{scp} (Linux/Mac).

\subsubsection{Archivo Individual}

\begin{lstlisting}
pscp -pw Pixelstore34.5 ^
  "server/routes/chat.js" ^
  root@177.7.41.162:/root/pixel-backend/server/routes/chat.js
\end{lstlisting}

\subsubsection{Todo el Backend}

\begin{lstlisting}
pscp -pw Pixelstore34.5 ^
  -r "server/*" ^
  root@177.7.41.162:/root/pixel-backend/server/
\end{lstlisting}

\subsubsection{Archivos Críticos del Backend}

\begin{itemize}
  \item \texttt{server/index.js} --- Entry point del servidor
  \item \texttt{server/database.js} --- Gestión de bases de datos JSON
  \item \texttt{server/services/discordService.js} --- Embeds de Discord
  \item \texttt{server/routes/orders.js} --- Lógica de pedidos
  \item \texttt{server/routes/chat.js} --- Mensajería en tiempo real
  \item \texttt{server/routes/reviews.js} --- Reseñas
  \item \texttt{server/routes/admin.js} --- Panel de administración
\end{itemize}

\subsection{Reiniciar el Servidor}

\begin{lstlisting}
plink -ssh -pw Pixelstore34.5 ^
  root@177.7.41.162 ^
  "cd /root/pixel-backend/server && pm2 restart pixel-server"
\end{lstlisting}

\subsection{Ver Logs}

\begin{lstlisting}
plink -ssh -pw Pixelstore34.5 ^
  root@177.7.41.162 ^
  "pm2 logs pixel-server"
\end{lstlisting}

\newpage

\section{Despliegue del Frontend (Build Estático)}

\subsection{Compilar el Frontend}

En la máquina local (Windows):

\begin{lstlisting}
cd C:\items\items\pixel_store_home_page_58ruox_dualiteproject
npm run build
\end{lstlisting}

Esto genera la carpeta \texttt{dist/}.

\subsection{Subir el Build al VPS}

\begin{lstlisting}
pscp -pw Pixelstore34.5 ^
  -r "dist/*" ^
  root@177.7.41.162:/var/www/frontend/
\end{lstlisting}

\subsection{Verificar la Instalación}

El contenido de \texttt{/var/www/frontend/} debe servirse via Nginx. Comprobar:

\begin{lstlisting}
plink -ssh -pw Pixelstore34.5 ^
  root@177.7.41.162 ^
  "ls -la /var/www/frontend/ && nginx -t && systemctl reload nginx"
\end{lstlisting}

\newpage

\section{Arquitectura del Sistema}

\subsection{Flujo de una Compra}

\begin{enumerate}
  \item Usuario selecciona productos en el frontend y paga.
  \item El frontend llama a \texttt{POST /api/orders} para crear el pedido.
  \item El backend guarda la orden en \texttt{server/db/orders.json}.
  \item Se emite \texttt{order-updated} via Socket.io al admin.
  \item El admin atiende el pedido desde el panel (\texttt{Admin.tsx}).
  \item Si el admin necesita contactar al cliente, usa \texttt{handleContactClient()} que crea un chat.
  \item Los mensajes viajan por \texttt{POST /api/chat/message} y se emiten en tiempo real por Socket.io.
  \item Cuando el pedido se completa, el usuario puede dejar una reseña.
  \item La reseña se envía a Discord mediante \texttt{discordService.js}.
\end{enumerate}

\subsection{Flujo de Mensajería en Tiempo Real}

\begin{enumerate}
  \item El cliente se conecta a Socket.io al cargar la página.
  \item Al abrir una orden, emite \texttt{join-order(orderId)} para unirse a la sala.
  \item Al abrir un chat, emite \texttt{join-chat(chatId)} para unirse a la sala del chat.
  \item Cuando alguien envía un mensaje, el backend emite \texttt{new-message} a ambas salas (\texttt{order-\{id\}} y \texttt{chat-\{id\}}).
  \item Esto asegura que el usuario reciba el mensaje aunque no haya abierto el chat explícitamente.
\end{enumerate}

\newpage

\section{Bases de Datos (JSON con lowdb)}

Todas las bases están en \texttt{server/db/}.

\begin{longtable}{|l|l|}
\hline
\textbf{Archivo} & \textbf{Contenido} \\ \hline
\endfirsthead
\hline
\textbf{Archivo} & \textbf{Contenido} \\ \hline
\endhead
\texttt{products.json} & Catálogo de productos \\ \hline
\texttt{orders.json} & Pedidos realizados \\ \hline
\texttt{chats.json} & Conversaciones y mensajes \\ \hline
\texttt{admins.json} & Cuentas de administradores \\ \hline
\texttt{users.json} & Usuarios registrados \\ \hline
\texttt{settings.json} & Configuración global de la tienda \\ \hline
\texttt{reviews.json} & Reseñas de clientes \\ \hline
\texttt{limiteds.json} & Items limiteds de Roblox \\ \hline
\texttt{mm2.json} & Items de Murder Mystery 2 \\ \hline
\texttt{coupons.json} & Cupones de descuento \\ \hline
\texttt{fortnite\_orders.json} & Pedidos de Fortnite \\ \hline
\texttt{group\_joins.json} & Solicitudes de unión a grupos \\ \hline
\end{longtable}

\newpage

\section{Discord Embeds}

Archivo: \texttt{server/services/discordService.js}

\subsection{Funciones Principales}

\begin{itemize}
  \item \texttt{sendOrderEmbed(order, type)} --- Embed de compra realizada.
    \begin{itemize}
      \item Colores: ingame = rojo, MM2 = rojo, robux = verde, limiteds = amarillo, Fortnite = cian.
      \item MM2 muestra \textit{Via Trade}, limiteds muestra \textit{Via Trade}, ingame/robux/Fortnite muestran \textit{Via Regalo}.
      \item Para ingame no se muestran nombres de productos, solo cantidad + total en robux.
    \end{itemize}
  \item \texttt{sendReviewEmbed(review)} --- Embed de reseña con imagen grande abajo.
  \item \texttt{getRobuxTotal(cart)} --- Calcula el total de robux de un carrito (busca en \texttt{products.json} por ID).
\end{itemize}

\subsection{Canales de Discord}

\begin{lstlisting}
Shop logs (compras):   canal configurado en discordService.js
Reviews:               canal configurado en discordService.js
Notificaciones admin:  via Socket.io (no Discord)
\end{lstlisting}

\newpage

\section{Comandos Rápidos}

\subsection{Despliegue Completo (Backend + Frontend)}

\begin{lstlisting}
:: 1. Compilar frontend
cd C:\items\items\pixel_store_home_page_58ruox_dualiteproject
npm run build

:: 2. Subir backend
pscp -pw Pixelstore34.5 -r "server/*" ^
  root@177.7.41.162:/root/pixel-backend/server/

:: 3. Subir frontend build
pscp -pw Pixelstore34.5 -r "dist/*" ^
  root@177.7.41.162:/var/www/frontend/

:: 4. Reiniciar servidor
plink -ssh -pw Pixelstore34.5 ^
  root@177.7.41.162 ^
  "cd /root/pixel-backend/server && pm2 restart pixel-server"
\end{lstlisting}

\subsection{Archivo Individual (Backend)}

\begin{lstlisting}
pscp -pw Pixelstore34.5 ^
  "server/routes/chat.js" ^
  root@177.7.41.162:/root/pixel-backend/server/routes/chat.js
plink -ssh -pw Pixelstore34.5 ^
  root@177.7.41.162 ^
  "cd /root/pixel-backend/server && pm2 restart pixel-server"
\end{lstlisting}

\newpage

\section{Solución de Problemas}

\subsection{El build falla}

\begin{lstlisting}
npx tsc --noEmit   # Verificar errores de TypeScript
npm run build      # Reintentar build
\end{lstlisting}

\subsection{El servidor no responde}

\begin{lstlisting}
plink -ssh -pw Pixelstore34.5 ^
  root@177.7.41.162 ^
  "pm2 status && pm2 logs pixel-server --lines 20"
\end{lstlisting}

\subsection{Los mensajes en tiempo real no llegan}

Verificar que el cliente está en la sala correcta:

\begin{lstlisting}
// En el navegador (consola):
socket.emit('join-order', 'ORDER_ID');
socket.emit('join-chat', 'CHAT_ID');
\end{lstlisting}

El backend emite \texttt{new-message} a las salas \texttt{chat-\{id\}} y \texttt{order-\{id\}}. Si el usuario no está en ninguna, no recibe el mensaje.

\subsection{El bot de Discord no envía mensajes}

\begin{lstlisting}
plink -ssh -pw Pixelstore34.5 ^
  root@177.7.41.162 ^
  "pm2 logs pixelito-bot --lines 30"
\end{lstlisting}

\newpage

\section*{Historial de Cambios}

\begin{tabular}{|l|l|l|}
\hline
\textbf{Fecha} & \textbf{Cambio} & \textbf{Responsable} \\ \hline
08/07/2026 & Versión inicial de la guía & Sistema \\ \hline
\end{tabular}

\end{document}
